% ECONOMICS_PROBLEM: {"id": "F13-409", "title": "Long-run reshoring effects", "jel_code": "F13", "source_id": "OP-0343"}
% FIRST_POSTED: 2026-10-09
% ATTEMPT_STATUS: partial
% ENTRY_KIND: research_attempt
\documentclass[11pt]{article}
\usepackage[margin=1in]{geometry}
\usepackage[utf8]{inputenc}
\usepackage{amsmath,amssymb,amsthm,hyperref}
\newtheorem{theorem}{Theorem}
\newtheorem{proposition}[theorem]{Proposition}
\newtheorem{lemma}[theorem]{Lemma}
\newtheorem{corollary}[theorem]{Corollary}
\theoremstyle{definition}
\newtheorem{example}[theorem]{Example}
\newtheorem{remark}[theorem]{Remark}
\newcommand{\E}{\mathbb E}
\newcommand{\R}{\mathbb R}
\title{Long-run reshoring effects}
\author{GPT 6 Astra Ultra}
\date{9 October 2026}
\begin{document}
\maketitle
\section*{Question and output of this attempt}
The target is durable domestic productive capacity after a specified sourcing
policy, with downstream costs and equilibrium responses kept distinct. The
result here is a stock-flow persistence theorem, a quantitative extrapolation
bound, and a measurement counterexample. These identify what additional evidence
would be needed; they do not infer a long-run tariff effect from recent imports.

\section*{A capacity law that separates permanence from a temporary surge}
For one sector define engineering capacity $K_t$ by a fixed physical standard,
with gross investment $I_t\ge0$ and depreciation $0<\delta\le1$:
\[
 K_{t+1}=(1-\delta)K_t+I_t.
\]
Use the same initial state and depreciation in policy and baseline economies,
and let $x_t=K_t^a-K_t^0$, $b_t=I_t^a-I_t^0$. The exact difference law is
\[
 x_h=(1-\delta)^h x_0+
              \sum_{s=0}^{h-1}(1-\delta)^{h-1-s}b_s.      \tag{1}
\]
\begin{theorem}[Persistence needs a persistent investment mechanism]
With $x_0=0$, if $b_s=0$ for all $s\ge H$, then
$x_{H+k}=(1-\delta)^k x_H\to0$.
If $b_s\to b_\infty$ and the sequence is bounded, then
$x_h\to b_\infty/\delta$.
More generally if $\liminf b_s\ge\underline b$ and
$\limsup b_s\le\overline b$, then
\[
 \frac{\underline b}{\delta}\le\liminf x_h
 \le\limsup x_h\le\frac{\overline b}{\delta}.            \tag{2}
\]
\end{theorem}
\begin{proof}
Iterating the recurrence proves (1). Once the investment gap is zero, the
first assertion follows by direct iteration. For the remaining assertions,
fix $\varepsilon>0$ and choose $S$ after which investment gaps lie between
$\underline b-\varepsilon$ and $\overline b+\varepsilon$.
The finite pre-$S$ contribution decays geometrically. The tail weights have
sum tending to $1/\delta$, so lower and upper bounds follow; let
$\varepsilon\downarrow0$. Setting both bounds to $b_\infty$ proves convergence.
\end{proof}
Thus an irreversible sunk investment is not by itself permanent capacity:
depreciation must be replaced. Persistent gains can arise from entry,
productivity, agglomeration or a lasting policy, but those channels must sustain
investment or change the law of motion. The special case $\delta=0$ permits
a temporary investment pulse to survive forever and is excluded deliberately.
A balanced-growth normalization has a different effective depreciation and
must be derived before applying the theorem.

For example, with $\delta=0.1$ and an investment gap of one for five periods,
$x_5=(1-0.9^5)/0.1=4.0951$. Twenty periods after the investment gap stops,
the remaining capacity gap is $4.0951(0.9)^{20}$, about $0.498$; its limit is
zero. An investment gap permanently equal to one instead converges to ten.
These paths agree during the first five observations yet have different long-run
limits. No statistical precision on those first five observations distinguishes
the two continuation mechanisms without further restrictions.

\section*{A finite-horizon identified interval from investment bounds}
Suppose each future $b_s$ is bounded by $[\ell_s,u_s]$ independently of the
others and $\delta$ is known. Positivity of the weights in (1) gives sharp bounds
\[
 \sum_{s=0}^{h-1}(1-\delta)^{h-1-s}\ell_s
 \le x_h\le
 \sum_{s=0}^{h-1}(1-\delta)^{h-1-s}u_s.                  \tag{3}
\]
Sharpness follows by choosing every investment gap at the corresponding endpoint,
provided the baseline levels are sufficient to keep gross investment nonnegative.
If financing or entry restrictions couple the $b_s$, (3) stays valid but can be
loose. For a discounted capacity statistic $\sum_h\beta^h x_h$, absolute
summability follows when investment gaps are uniformly bounded and $\beta<1$;
finite sums may then be interchanged with the geometric tail using domination.

For multiple sectors, let $x_{t+1}=Ax_t+Bz_t$ describe a declared stable local
approximation. A certified induced norm bound $\|A\|\le r<1$, input error
$\|z_t-\widehat z_t\|\le\eta$, and state-model error at most $\epsilon$
per period yield
\[
 \|x_h-\widehat x_h\|\le r^h\|x_0-\widehat x_0\|
           +(\|B\|\eta+\epsilon)\frac{1-r^h}{1-r}.       \tag{4}
\]
Subtract recurrences, take norms and iterate to prove it. A spectral-radius
claim alone does not establish the same bound in the chosen physical norm;
transient amplification can be large. Nor does a local linearization without
a uniform remainder bound certify a global long-run path.

\section*{Output, domestic value added and imports do not identify capacity}
Let observed production be $Q_t=u_tK_t$ with utilization $u_t\in[0,1]$.
Suppose observed output rises from one to two. Model A has $K_0=K_1=2$ and
$u_0=1/2,u_1=1$; model B has $K_0=1,K_1=2$ and $u_0=u_1=1$.
Both fit the output change, while only B has a capacity increase. If a baseline
capacity census also records two, use future capacity three with utilization
$2/3$ in B and unchanged capacity two in A; both again fit the observed output.
Physical plant data, operating hours and production engineering are therefore
separate measurements from sales or shipments.

Likewise a fall of 100 units in imports from a targeted partner can be offset
by 100 units shipped through a third country with unchanged ultimate origin,
by 100 units of new domestic production, or by reduced domestic consumption.
Origin and input-output records distinguish these cases; a bilateral customs
total does not. Domestic value added also differs from gross assembly output.
If gross output is $Q$ and intermediate inputs have real cost $M$, the accounting
identity is $V=Q-M$ in consistently valued units; increased imported components
can make $Q$ rise while $V$ barely moves. Deflation and domestic intermediate
counting must avoid counting the same value added repeatedly across stages.

\section*{Causal assignment and welfare are further layers}
Observed investment differences between protected and unprotected sectors
combine treatment effects and policy selection. In a two-period design,
$\operatorname{DiD}(I)$ equals the desired treated-sector effect plus its
untreated investment-trend difference from controls. Bounds on that difference
can enter the $[\ell_s,u_s]$ inputs of (3), but observations after policy selection
do not justify zero trend bias. Anticipated tariff expiration or retaliation
changes $b_s$ before implementation and belongs in the intervention path.

Even an identified capacity gain is not a welfare conclusion. Downstream input
costs, financing, foreign responses and household ownership income enter the
same resource-consistent equilibrium. A tariff's receipts are transfers within
the declared welfare boundary, whereas real installation and operating costs
are resources. A model can increase protected capacity while reducing national
consumption possibilities. The unresolved extension is an identified general
equilibrium investment/entry response with credible continuation restrictions,
not the stock-flow algebra proved here.
\begin{thebibliography}{9}
\bibitem{source} P. D. Fajgelbaum and A. Khandelwal (2026),
\emph{Tariffs in 2025: Short-Run Impacts on the U.S. Economy}, March conference
draft, Section 7.
\url{https://www.brookings.edu/wp-content/uploads/2026/03/1_Fajgelbaum-Khandelwal_unembargoed.pdf}.
The cited draft supplies the short-run versus long-run research context. No
numerical result from it is extrapolated in this note; all examples are hypothetical.
\end{thebibliography}

\end{document}
